% Sucesor del articulo V: incorporacion focal F051.
% Fragmento bilingue emparejado con la homologa EN; no compilar aisladamente.
\subsection{Observateur premier des vacances et résolution modale}
\label{sec:f051-prime-vacancy-observer}

Dans la chaîne causale de l’état HMT déjà engendré, le produit nonadique, son
coproduit multiplicatif réduit et la projection spectrale
\(P_{\rm irr}\) produisent les formes irréductibles ; l’évaluation
\(\operatorname{ev}_9\) les reconnaît ensuite comme nombres premiers ordinaires. La
frontière \(729/1000\) produit le calendrier des vacances. L’observateur
\(\mathcal C_{\varepsilon,X}\) compose ces deux objets internes avec l’horloge
logarithmique. Le crible conventionnel intervient uniquement comme lecteur
externe de la reproduction numérique.

\subsubsection{Calendrier et équivalence des deux écritures}

Soit
\begin{equation}
 \varepsilon_{\rm vac}
 =1-\log_{10}9
 =\log_{10}\!\left(\frac{10}{9}\right),
 \qquad
 z_n=\left\lceil(n+1)\varepsilon_{\rm vac}\right\rceil
      -\left\lceil n\varepsilon_{\rm vac}\right\rceil ,
 \quad n\geq1.
 \label{eq:f051-epsilon-calendar}
\end{equation}
La quantité \(\varepsilon_{\rm vac}\) est irrationnelle. En effet, si
\(\varepsilon_{\rm vac}=a/b\in\mathbb Q\), alors
\((10/9)^b=10^a\), ce qui contredit l’unicité de la factorisation
en nombres premiers par la présence d’une puissance non nulle de \(3\). Par conséquent,
ni \(n\varepsilon_{\rm vac}\) ni
\((n+1)\varepsilon_{\rm vac}\) ne sont des entiers pour \(n\geq1\), et
\begin{equation}
 z_n=\left\lfloor(n+1)\varepsilon_{\rm vac}\right\rfloor
      -\left\lfloor n\varepsilon_{\rm vac}\right\rfloor
 \in\{0,1\}.
 \label{eq:f051-ceil-floor}
\end{equation}
La restriction \(n\geq1\) est nécessaire à cette équivalence : en
\(n=0\), la convention avec plafonds et la convention avec planchers ne coïncident pas.

\subsubsection{Chaîne typée et observateur de corrélation}

Pour \(X\geq2\), la chaîne causale peut s’écrire sans perte de types
comme
\begin{equation}
 \operatorname{Ran}(P_{\rm irr})
 \xrightarrow{\ \operatorname{ev}_9\ }
 \mathbb P
 \xrightarrow{\ \iota\ }
 \{0,1\}\times\mathbb R_{>0},
 \qquad
 \bigl((z_p,\log p)\bigr)_{p\in\mathbb P_{\leq X}}
 \xrightarrow{\ \mathcal C_{\varepsilon,X}\ }
 \mathbb R,
 \label{eq:f051-typed-chain}
\end{equation}
où
\[
 \iota(p)=(z_p,\log p),
 \qquad
 \mathcal C_{\varepsilon,X}\colon
 \prod_{p\in\mathbb P_{\leq X}}
       (\{0,1\}\times\mathbb R_{>0})\longrightarrow\mathbb R,
 \qquad
 \mathcal C_{\varepsilon,X}\bigl((a_p,\ell_p)_{p\leq X}\bigr)
 =\sum_{p\leq X}(a_p-\varepsilon_{\rm vac})\ell_p .
\]
Par conséquent,
\begin{equation}
 \theta_{\rm vac}(X)=\sum_{p\leq X}z_p\log p,
 \qquad
 \theta(X)=\sum_{p\leq X}\log p,
 \qquad
 R_{\rm vac}(X)=\theta_{\rm vac}(X)
                   -\varepsilon_{\rm vac}\theta(X).
 \label{eq:f051-rvac-chebyshev}
\end{equation}
Le centrage exact de l’observateur est
\(\varepsilon_{\rm vac}\theta(X)\) ; l’expression
\(\varepsilon_{\rm vac}X\) correspond à une autre normalisation. Deux préfixes
avec le même nombre de vacances peuvent produire des valeurs différentes de
\(R_{\rm vac}\), car l’observateur conserve la position première et son poids
logarithmique.

\subsubsection{Lecteurs de Dirichlet et domaines analytiques}

L’observateur fini et les deux lecteurs de Dirichlet constituent
des réalisations distinctes et compatibles :
\begin{align}
 Z^{+}_{\rm vac}(s)
   &=\sum_{n\geq1}\frac{z_n}{n^s}
     =\varepsilon_{\rm vac}\zeta(s)+H_{\rm vac}(s),
 & H_{\rm vac}&\in\mathcal O(\operatorname{Re}s>0),
 \nonumber\\
 Z^{\times}_{\rm vac}(s)
   &=\prod_p(1-p^{-s})^{-z_p}
     =\zeta(s)^{\varepsilon_{\rm vac}}G_{\rm vac}(s),
 & \operatorname{Re}s&>1,
 \label{eq:f051-dirichlet-readers}\\
 \log G_{\rm vac}(s)
   &=\sum_p(z_p-\varepsilon_{\rm vac})
       \log(1-p^{-s})^{-1},
 & \operatorname{Re}s&>1.
 \nonumber
\end{align}
La première série et le produit d’Euler convergent absolument pour
\(\operatorname{Re}s>1\). La discrépance partielle du mot mécanique
est bornée, et la sommation d’Abel prolonge \(H_{\rm vac}\) de manière
holomorphe à \(\operatorname{Re}s>0\). La puissance
\(\zeta(s)^{\varepsilon_{\rm vac}}\) est définie ici au moyen du logarithme
d’Euler dans le demi-plan de convergence absolue.

\subsubsection{Résolution de Fourier et réalité de l’observateur}

Définissons, pour \(k\in\mathbb Z\setminus\{0\}\),
\begin{equation}
 c_k=\frac{e^{2\pi i k\varepsilon_{\rm vac}}-1}{2\pi i k},
 \qquad
 S_k(X)=\sum_{p\leq X}(\log p)
             e^{2\pi i k p\varepsilon_{\rm vac}}.
 \label{eq:f051-fourier-data}
\end{equation}

\begin{proposition}[Décomposition modale de la corrélation première]
\label{prop:f051-fourier-resolution}
Pour tout \(X\geq2\),
\begin{equation}
 R_{\rm vac}(X)
 =\lim_{K\to\infty}
   \sum_{0<|k|\leq K}c_kS_k(X)
 =\lim_{K\to\infty}
   \sum_{0<|k|\leq K}
   \left(1-\frac{|k|}{K+1}\right)c_kS_k(X).
 \label{eq:f051-fourier-pair}
\end{equation}
La première limite est définie au moyen de sommes partielles symétriques et la
seconde au moyen de moyennes de Fejér ; ce sont exactement les deux notions
de convergence affirmées. La convergence absolue en \(k\) reste en dehors
des hypothèses et de la conclusion.
\end{proposition}

\begin{proof}
Soit
\[
 f_{\varepsilon}(x)
 =\mathbf1_{(1-\varepsilon_{\rm vac},1)}(\{x\})
  -\varepsilon_{\rm vac}.
\]
Ses coefficients de Fourier non nuls sont précisément \(c_k\).
L’irrationalité de \(\varepsilon_{\rm vac}\) empêche que
\(p\varepsilon_{\rm vac}\) atteigne une extrémité de l’intervalle ; par
\eqref{eq:f051-ceil-floor},
\(f_{\varepsilon}(p\varepsilon_{\rm vac})
 =z_p-\varepsilon_{\rm vac}\). Le théorème de Dirichlet donne la convergence
des sommes symétriques en ces points de continuité, et le théorème de
Fejér donne la convergence des moyennes. La finitude de l’ensemble
\(p\leq X\) permet d’intervertir la limite et la somme première pondérée aux
sens symétrique et de Fejér. L’estimation
\(|c_k|=O(|k|^{-1})\) n’implique pas la sommabilité absolue et ne
permet donc pas de réordonnement absolu.

Enfin,
\(c_{-k}=\overline{c_k}\) et
\(S_{-k}(X)=\overline{S_k(X)}\). Les modes \(k\) et \(-k\) contribuent
\(2\operatorname{Re}(c_kS_k(X))\) ; on en retrouve directement la
réalité de \(R_{\rm vac}(X)\).
\end{proof}

\subsubsection{Témoins numériques et reproduction indépendante}

La Table~\ref{tab:f051-rvac-powers} présente les six coupures finies
disponibles. Les valeurs de \(R_{\rm vac}\) reproduisent tous les chiffres
enregistrés et la colonne normalisée conserve l’arrondi de la source
numérique.

\begin{table}[htbp]
\centering
\small
\begin{tabular}{rrrrr}
\toprule
\(X\) & \(\pi(X)\) & \(\pi_{\rm vac}(X)\) & \(R_{\rm vac}(X)\)
 & \(R_{\rm vac}(X)/\sqrt X\)\\
\midrule
\(10^3\) & 168 & 7 & \(-4.553590011317908\) & \(-0.143997159663565\)\\
\(10^4\) & 1\,229 & 61 & \(41.334631595982586\) & \(0.413346315959826\)\\
\(10^5\) & 9\,592 & 427 & \(-141.1579619980301\) & \(-0.446380669781268\)\\
\(10^6\) & 78\,498 & 3\,595 & \(32.21713602110386\) & \(0.032217136021104\)\\
\(10^7\) & 664\,579 & 30\,384 & \(-418.47581825715594\) & \(-0.132333673139529\)\\
\(10^8\) & 5\,761\,455 & 263\,841 & \(4021.0180144347314\) & \(0.402101801443473\)\\
\bottomrule
\end{tabular}
\caption{Observateur des vacances sur les nombres premiers dans six coupures finies.}
\label{tab:f051-rvac-powers}
\end{table}

La table complète des trente modes se trouve dans
\texttt{S\_k\_1\_30\_X1e8.csv}. Ses cinq plus grandes valeurs de
\(|S_k(10^8)|\), ordonnées par module, sont :

\begin{table}[htbp]
\centering
\scriptsize
\setlength{\tabcolsep}{3.5pt}
\begin{tabular}{rrrrrr}
\toprule
\(k\) & \(\operatorname{Re}S_k\) & \(\operatorname{Im}S_k\)
 & \(|S_k|\) & \(|S_k|/\sqrt X\) & \(|S_k|/(\sqrt X\log X)\)\\
\midrule
19 & 30646.583055876723 & -39786.84285255695 & 50221.56824686792 & 5.0221568246867925 & 0.2726368745267788\\
24 & -1325.0060816425494 & 47631.087438932576 & 47649.51344695589 & 4.764951344695589 & 0.25867400944234675\\
21 & 1214.7582244307505 & -44824.22356680968 & 44840.68081453668 & 4.484068081453668 & 0.24342575303172864\\
26 & -26879.47488364538 & 33364.6769815865 & 42845.161221614406 & 4.284516122161441 & 0.23259271368502904\\
1 & 40391.94841735065 & 13384.21089257329 & 42551.693246765084 & 4.255169324676508 & 0.23099956965887428\\
\bottomrule
\end{tabular}
\caption{Cinq plus grands modes dans l’échantillon \(1\leq k\leq30\) pour
\(X=10^8\). Dans ce domaine d’échantillonnage,
\(\lvert S_{19}(10^8)\rvert\) est le maximum.}
\label{tab:f051-top-five-modes}
\end{table}

Le calcul indépendant des quatre premières lignes, obtenu à partir de
\eqref{eq:f051-epsilon-calendar} avec un crible utilisé comme lecteur de
vérification, reproduit les valeurs jusqu’à \(10^6\) avec des erreurs absolues
inférieures à \(2.1\times10^{-11}\). Les lignes de \(10^7\) et \(10^8\) et les
trente modes sont des témoins numériques finis du calcul historique. La
projection \(P_{\rm irr}\) de \eqref{eq:f051-typed-chain} conserve le rôle
de générateur interne.

\subsubsection{Condition analytique pour la borne critique}

L’estimation
\begin{equation}
 R_{\rm vac}(X)=O\!\left(X^{1/2}\log^B X\right)
 \label{eq:f051-conditional-bound}
\end{equation}
contrôlerait le terme premier centré du facteur \(G_{\rm vac}\). Sa
démonstration exige des bornes pour \(S_k(X)\) uniformes en \(k\) et compatibles
avec la somme symétrique ou de Fejér. Les tables caractérisent une fenêtre
finie et n’apportent pas cette uniformité. Le résultat démontré est
l’identité modale exacte pour chaque \(X\) fini ; la borne asymptotique critique
reste conditionnée aux estimations modales uniformes.
